% ============================================================
% motivation.tex —— 动机与历史回顾（Motivation and Historical Perspective）
% 底本：Springer LNM 2386，PDF p019–p040，印刷 xxi–xlii
% 页锚点：正文中以「原文 p.xxi」形式的注释标记页界（本书前言用罗马数字页码）
% 本块无编号公式；插图仅 Fig 0.1
% ============================================================

\chapter*{动机与历史回顾}

% [原文 p.xxi]

\ChapterAbstract{我们在此的目标是给出一个简短的——因而必然是不完整的——编年式概述，梳理来自高能物理与泰希米勒理论两方面的关键思想，这些思想共同促成了谱网络的诞生。我们将着重说明这两个领域如何相互充实：一方的概念常常为另一方开辟全新的研究方向，有时它们各自独立发展，但晚近则通过谱网络这一框架汇聚起来。我们首先讨论量子场论的重要方面——这是描述量子相对论性现象的框架——重点是 Yang--Mills 理论的引入与瞬子的发现。接下来转向数学中的规范理论，它正是从 Yang--Mills 理论中瞬子的研究里生长出来的。我们还会说明高阶泰希米勒理论如何从规范理论的研究中演化出来，并显著推广了经典的泰希米勒空间概念。第三节考察超对称量子场论的若干方面，主要把它当作一个富有成效的工作假设，使我们得以探索量子场论中原本可能无法触及的深层内容。我们将特别强调电磁对偶——它既促成了对具有扩展超对称的四维量子场论中 BPS 态的研究，也促成了 Seiberg--Witten 理论的发展。最后，我们讨论弦论技巧如何拓宽了可用 Seiberg--Witten 分析处理的理论类，并使研究这些理论的方法系统化。这尤其带来了一大类被称为 S 类理论的理论，其中 Hitchin 系统居于中心地位。在这些理论中对 BPS 态的研究，揭示了与 Hitchin 分支（高阶泰希米勒空间的重要例子）理论之间出人意料的联系，并最终促成了谱网络的发展；而谱网络在高阶泰希米勒空间的研究中已迅速引起广泛关注。}

% [原文 p.xxii]

\section*{量子场论与规范理论}

量子理论是描述小尺度（通常是原子尺度或更小）自然行为的基本框架。它在 20 世纪初的几十年间逐步发展起来。量子理论的一个核心方面是：量子系统的态由可分 Hilbert 空间中的射线表示，这一数学构造称为态 Hilbert 空间 $(\mathcal{H},\langle\cdot|\cdot\rangle)$。量子理论的发展与 Hilbert 空间理论的演进密切相关，John von Neumann 等著名科学家在其中作出了重要贡献。Wigner 1931 年的定理是量子理论的重要基石；它断言量子系统的任何对称性都必须对应于 $\mathcal{H}$ 的酉或反酉、线性或反线性变换。

狭义相对论是描述空间与时间关系的理论，它建立在两条核心原则之上：相对性原理，以及经实验验证的假设——真空中的光速在所有惯性参考系中都相同。该理论由 Einstein 于 1905 年提出，但建立在此前众多科学家的贡献之上。在狭义相对论中，空间与时间被统一为称为时空的单一框架，它在数学上被建模为带符号差 $(1,3)$ 度量的四维伪黎曼微分流形。例如，平坦时空由 Minkowski 时空表示，即装备伪黎曼度量 $\operatorname{Diag}(-1,1,1,1)$ 的 $\R^4$。狭义相对论的发展深受 19 世纪几何学进展的影响，尤其是 Riemann 的开创性贡献。

\subsection*{量子场论}

在非相对论量子力学中，孤立量子系统的态可以用波函数描述。波函数模的平方给出概率密度，从而赋予系统的物理量以概率解释。波函数的演化由 Schrödinger 方程支配，该方程不对称地对待时间与空间，因而与狭义相对论的框架不相容。Schrödinger 方程存在相对论性推广，例如 Klein--Gordon 方程与 Dirac 方程。然而，这些相对论性方程的解缺乏直接的概率解释。归根结底，这是因为在狭义相对论中粒子数不守恒（例如 [19] 中对此有讨论）。

这一认识促成了量子场论（QFT）的发展，它为量子相对论系统提供了框架。特别地，Minkowski 时空中的一切量子场论都把时空的对称性——具体说是 Poincaré 群 $\R^4\rtimes \mathrm{O}(3,1)$——作为其对称性的基本子群加以尊重。

量子场论的发展在很大程度上得益于对光与物质之间量子相互作用的研究，特别是 20 世纪 40 年代发展起来的量子电动力学（QED）。QED 是经典 Maxwell 理论的量子场论推广，它把其激发描述电子与正电子的荷电物质场纳入框架。
% [原文 p.xxiii]
它作为规范理论的最简单例子出现，其中规范群为 $\mathrm{U}(1)$。QED 的成功，尤其是在精确预言实验结果方面的成功，使它成为量子场论更广泛发展中的基石，并为粒子物理中更复杂的规范理论提供了原型。

\subsection*{重整化}

在量子电动力学（QED）中，许多计算会给出无穷大，这要求 20 世纪 40 年代末广泛发展起来的正则化方法。然而人们逐渐认识到，这些无穷大并不是需要消除的问题，而是量子场论的内在特征。这一认识导致了重整化（renormalization）的概念：在量子场论中，关联函数依赖于进行实验的能量标度。在 Wilson 的范式下，这使理论的参数（如质量、电荷与耦合常数）产生对能量的依赖。例如，无量纲的精细结构常数 $\alpha$ 通过
\[
4\pi\epsilon_0\hbar c\,\alpha(\mu)=e^2(\mu)
\]
与基本电荷 $e$ 相联系，它随 $\mu$ 变化而呈现能标依赖。低能下测得它约为 $1/137$，而在 $Z$ 玻色子的能标 $m_Z$（约 90 GeV）处增大到约 $1/127$。这种变化之所以出现，是因为每个电子都与它自己产生的、依赖标度的电磁场相互作用。

重整化理论在凝聚态物理概念的显著推动下大大增强，尤其促成了 20 世纪 60、70 年代 Wilson 重整化群的发展。这一框架描述理论的参数 $g_1,\dots,g_k$ 如何随能标 $\mu$ 以在 $\mu$ 中局部的方式演化。具体地，可以把这种变化写成
\[
\mu\frac{\dd g_i}{\dd\mu}(\mu_0)=:\beta_{g_i}\bigl(g_1(\mu_0),\dots,g_k(\mu_0)\bigr),\qquad \text{对一切 } i=1,\dots,k,
\]
其中 $\beta$ 是依赖于理论、在能标 $\mu_0$ 处取值的参数函数。这一关系在理论的参数空间中诱导出一个流，称为重整化群流（renormalization group flow）。在此框架下，参数空间中的某些点可能在该流下固定；这些点对应的理论具有标度不变性。这一现象的一个突出例子出现在共形场论的语境中：对共形场论而言，时空的共形对称群是其对称性的子群，从而特别蕴含标度不变性。

% [原文 p.xxiv]

\subsection*{Yang--Mills 理论}

量子场论框架的一个基石是 1954 年 Yang--Mills 理论的引入 [372]。这些理论推广了 Maxwell 理论；二者虽然都是规范理论，但 Yang--Mills 理论可以与任意约化紧单李群 $G$ 相联系，而不只是电磁学中所用的交换群 $\mathrm{U}(1)$。在此语境下，规范场充当 Maxwell 理论中四维势 $A_\mu$ 的推广，并被表示为时空上 $G$-主丛 $P$ 上的一个联络。除规范场外，还可以加入其他场，例如与主丛 $P$ 相伴的向量丛的截面。规范场的运动方程即所谓的 Yang--Mills 方程，它们构成推广 Maxwell 方程的一个偏微分方程组。这些理论成功地描述了原子核内部的弱相互作用与强相互作用。
此外，Yang--Mills 方程在数学的诸多分支中得到了广泛应用，稍后我们将对此加以讨论。

% [Tutor Note: 历史与文献线索：Yang--Mills 的原始论文是 C. N. Yang, R. L. Mills, "Conservation of Isotopic Spin and Isotopic Gauge Invariance", Phys. Rev. 96 (1954) 191，即本文 [372]。当时无质量的规范场与可观测短程力相矛盾，Pauli 曾在报告会上当场质询；规范场质量最终由 Higgs 机制给出（见后面"渐近自由"一节的电弱理论）。另：书中说 Yang--Mills 理论可联系"任意约化紧单李群"，严格说规范群通常取紧半单（或紧约化）李群——QCD 取 SU(3)，电弱理论取 SU(2)×U(1)，都是其特例。]

\subsection*{反常}

手征反常最早由 Adler 与 Bell--Jackiw 于 1969 年研究，其动因是中性 $\pi^0$ 介子衰变为两个光子 $\pi^0\to\gamma+\gamma$ 的反常速率。一般地，场论中的经典对称性若不能作为量子对称性保持下来，就称它是反常的（anomalous）。手征反常与 Atiyah--Singer 指标定理（1963）有深刻联系，这一联系尤以 Fujikawa 1984 年对手征反常的分析著称。此后，这一概念得到了显著拓宽。一个相关且同样基本的概念是整体对称性的 't Hooft 反常（1980），它们是内蕴的重整化群不变量的量。't Hooft 反常对量子场论的动力学施加严格约束，即使在缺乏标准计算工具时也是如此。

\subsection*{渐近自由、QCD 与禁闭}

对弱核力的研究，尤其是通过吴实验（1956），最终导致 1967 年 Weinberg--Salam 电弱理论的建立，它统一了电磁力与弱力。该理论是规范群为 $\SU(2)\times\mathrm{U}(1)$ 的 Yang--Mills 理论的特例，随着 2012 年 CERN 发现 Higgs 玻色子而获得完全的实验确认。

至于强核力，夸克模型于 20 世纪 60 年代初被提出，而夸克在 1968 年被实验观测到。强相互作用最好由 $\SU(3)$ Yang--Mills 理论描述这一认识于 1973 年形成，由此产生了今天所说的量子色动力学（QCD）。电弱理论与 QCD 共同构成粒子物理标准模型的关键组成部分，
% [原文 p.xxv]
而标准模型在 20 世纪 70 年代末已发展完备。QCD 的进展尤其得益于 1973 年渐近自由（asymptotic freedom）的发现。

QCD 是渐近自由的，意思是随能标升高，该理论越来越像自由量子场论，即没有相互作用、可以精确求解的理论。在远高于 $\Lambda_{QCD}\sim 330$ MeV 的能标处，QCD 可近似为夸克与胶子的自由理论。夸克–胶子相互作用的强度由规范耦合常数 $\alpha_s$ 决定，它随能标升高而减小，在极高能下趋于零。标度 $\Lambda_{QCD}$ 称为 QCD 的动力学标度，它由量纲嬗变（dimensional transmutation）而出现。

在弱耦合量子场论中，即相互作用几乎可忽略时——如极高能下的 QCD——可以使用多种微扰技巧进行分析，例如 Feynman 图。这些技巧把物理量表示为耦合常数的级数展开，使人们能计算前几项，并在耦合越弱时给出越来越精确的近似。理论越接近自由，这种微扰近似就越精确。

相反，当量子场论强耦合时，能直接分析它们的既有技巧极少。对能标低于 $\Lambda_{QCD}$ 的 QCD 尤其如此。不幸的是，正是在这一强耦合区域中出现了关键现象，例如夸克与胶子被禁闭（confinement）在质子与中子（它们构成原子核）之中。严格理解禁闭至今仍是 Clay 研究所千禧年大奖难题之一。

\subsection*{电磁对偶}

Maxwell 方程如下：
\begin{align*}
\nabla\cdot\mathbf{E}&=\frac{\rho}{\epsilon_0},\qquad \nabla\cdot\mathbf{B}=0,\\
\nabla\times\mathbf{E}&=-\frac{\partial\mathbf{B}}{\partial t},\qquad
\nabla\times\mathbf{B}=\mu_0\left(\mathbf{J}+\epsilon_0\frac{\partial\mathbf{E}}{\partial t}\right).
\end{align*}
在真空中没有电荷或电流（即 $\rho=0$、$\mathbf{J}=\mathbf{0}$），方程在变换 $(\mathbf{E},\mathbf{B})\to(\mathbf{B},-\mathbf{E})$ 下保持不变（假定采用 $\mu_0\epsilon_0=c^2=1$ 的自然单位）。这一变换称为电磁对偶（electric-magnetic duality）。为了在有源时恢复这一对偶性，需要同时引入磁荷密度 $\rho_m$ 与磁流 $\mathbf{J}_m$，它们由磁单极子产生。1931 年，Dirac 著名地证明了今天所谓的 Dirac 量子化条件，表明磁单极子的存在将要求电荷的量子化，
% [原文 p.xxvi]
而这一现象有重要的实验证据。

't Hooft [209] 与 Polyakov [305] 于 1974 年为磁单极子提供了进一步的理论支持，他们证明磁单极子作为某些非交换 Yang--Mills 理论中的经典孤子解自然出现。在他们的工作基础上，Julia 与 Zee 于 1975 年证明，同时携带电荷与磁荷的粒子（称为双子，dyon）在此语境中也作为孤子构型出现 [219]。此外，'t Hooft、Polyakov 以及 Julia--Zee 的分析可推广到取代 Maxwell 理论的非交换 Yang--Mills 理论。这一推广使 Goddard、Nyuts 与 Olive 于 1975 年证明：规范群为 $G$ 的 Yang--Mills 理论中的磁单极子属于所谓 GNO 对偶规范群 $G^\vee$ [165] 的表示，而 $G^\vee$ 也是 $G$ 的 Langlands 对偶（Langlands 纲领始于 Langlands 1967 年致 Weil 的信）。这一奠基性工作促使 Montonen 与 Olive 于 1977 年对（经典）Yang--Mills 理论表述了电磁对偶的非交换版本 [271]。

\subsection*{瞬子}

Yang--Mills 理论中的瞬子由 Belavin、Polyakov、Schwarz 与 Tyupkin 于 1975 年发现 [45]。形式上，瞬子是欧氏 Yang--Mills 方程的一个具有有限非零作用量的经典解。这意味着该解是一个对偶或自对偶的主 Yang--Mills 联络。物理上，瞬子捕捉了 Yang--Mills 理论中重要的非微扰效应，使人们能够计算不同经典真空之间的隧穿振幅。这反过来有助于理解这些理论中量子真空的结构。由于其非微扰本性，弱耦合时 Yang--Mills 理论中的瞬子效应极其微小，这解释了为何这些量子场论形式体系的稳健预言迄今尚未被直接观测到。
% [Tutor Note: "对偶或自对偶"指曲率满足 *F = ±F：自对偶取 +、反自对偶取 −；瞬子的作用量达到最小值 8π²/g²（BPST，1975）。ADHM（Atiyah--Drinfeld--Hitchin--Manin，1978）用四元数矩阵与 CP^3 上的全纯丛同构类参数化全部有限能量解，这正是文中"把问题翻译成 CP^3 上全纯向量丛"的含义。另：原文把 Nuyts 印作 "Nyuts"（p.xxvi），照印。]

\section*{自对偶方程与几何}

\subsection*{数学中的规范理论}

求解自对偶 Yang--Mills 方程的问题，对从事量子场论的物理学家变得极为重要，其中一些解在 Yang 与 Mills 的工作之后的几年里就已被构造出来。然而，求解这些方程的问题具有非常清晰的微分几何含义，并且由 Singer 于 1977 年带到牛津的几何学派。$\R^4$ 上主 $\SU(2)$-丛 $P$ 上的一个联络 $A$ 满足自对偶 Yang--Mills 方程，如果它的曲率 $F(A)$ 在 $\Omega^2(\R^4;\operatorname{ad}(P))$ 上的 Hodge 星算子下不变：

% [原文 p.xxvii]

\[
\star F(A)=F(A).
\]

Atiyah--Singer 指标定理所提供的解比物理学家已经产生的解更多。然而此时的目标是以具体形式找出所有解。为此，Atiyah 与 Ward [34] 首先改写了 Penrose 与 Ward [299], [353] 关于平坦时空的扭量几何与扭量运动学的工作，把复几何的思想引入数学物理。更精确的理解随 Atiyah--Drinfeld--Hitchin--Manin（ADHM）构造 [36] 的出现而来，它使得可以描述 4-球面 $S^4$ 上自对偶 Yang--Mills 方程的全部有限能量解。这一处理建立在扭量构造之上，后者把问题转化为关于 $\C P^3$ 上全纯向量丛的问题，从而把这些方程的研究与代数几何中的丛论联系起来。

下一个奠基性的步骤由 Atiyah 与 Bott 于 1983 年迈出 [33]，他们研究了任意亏格黎曼曲面上的 Yang--Mills 方程。在此情形下，Yang--Mills 泛函具有任意高 Morse 指标的临界点，并且存在一个完美的 Morse 理论，提供关于 Yang--Mills 极小点空间的信息。Narasimhan--Seshadri 关于用稳定全纯向量丛描述不可约酉联络的途径 [278] 似乎自然地契合这一图景，因为由基本群表示产生的丛恰好给出临界点。于是，一个关于稳定主丛与 Hermite--Einstein 联络存在性之间关系的更广阔图景开始浮现。是 Donaldson [115] 把这一问题重新表述为无穷维矩映射问题。大约同时，Hitchin--Karlhede--Lindstr\"om--Ro\v{c}ek [201] 给出了超凯勒流形的两种构造，使超对称与现代微分几何之间的联系更为透明。

这些工作指向一个具有极丰富几何解释的解空间的存在。确实，Hitchin 于 1987 年 [198] 通过解析地构造黎曼曲面上 $\SU(2)$ 自对偶 Yang--Mills 方程的解模去规范变换所得的模空间 $\mathcal{M}$（作为一个完备超凯勒流形），最终清楚地展示了这一点。在这一特定语境中，方程可写为
\[
F(A)+[\Phi,\Phi^*]=0,\qquad d''_A\Phi=0,
\]
其中 $A$ 是黎曼曲面上主丛 $P$ 上的联络，$\Phi$ 是黎曼曲面上取值于 $P$ 的复李代数丛的全纯 $(1,0)$-形式。

这些方程的解还与黎曼曲面的内蕴结构有密切关系。即，由映射 $(A,\Phi)\mapsto(A,-\Phi)$ 在 $\mathcal{M}$ 上诱导的对合的不动点集，有一个分量对应于 Euler 类等于 $2g-2$ 的平坦丛，其中 $g$ 表示底曲面的亏格。自对偶方程的解随后给出曲面上
% [原文 p.xxviii]
常负曲率 $-4$ 的度量，从而给出从某个解集到该曲面的泰希米勒空间的自然微分同胚。

\subsection*{经典泰希米勒理论}

对亏格 $g\ge 2$ 的闭连通定向拓扑曲面 $S$，$S$ 的泰希米勒空间 $\mathcal{T}(S)$ 定义为 $S$ 上所有带标记双曲结构模去同痕（即可连续形变为恒等映射的微分同胚）所得的空间。这个空间带有典范的复流形结构，并容许丰富的度量。对这些不同结构的几何特征的研究有着悠久传统，其开端包括 Riemann、Poincar\'e 与 Klein 等人。Poincar\'e 与 Koebe 于 1907 年建立的单值化定理也把 $\mathcal{T}(S)$ 描述为 $S$ 上复结构模去同痕于恒等映射的微分同胚所得的模空间。

对这一模空间的深入研究由 Teichm\"uller 于 20 世纪 30 年代末进行，此空间自此以他命名。Teichm\"uller 认识到，为研究这个空间，不仅考虑共形映射、而且考虑拟共形映射是有用的。拟共形映射此前已由 Gr\"otzsch 于 1928 年引入复分析；它们是平面区域之间的同胚，一阶地看把小圆变成有界偏心率的椭圆。在诸多结果中，Teichm\"uller 在 $\mathcal{T}(S)$ 上引入了一个度量（今称泰希米勒度量），并发现了极值拟共形映射与全纯二次微分之间的密切关系，从而断言 $\mathcal{T}(S)$ 同胚于 $\R^{6g-6}$。这一基本途径后来由 Ahlfors 与 Bers 在 20 世纪 50、60 年代用更解析的方法进一步推广，并为 Teichm\"uller 的结果提供了严格证明。

此后是 Thurston 为本课题引入了更具几何意味的风味，特别是通过他关于映射类群在 $\mathcal{T}(S)$ 上作用的研究。该群是 $S$ 的全部微分同胚模去同痕所得的群，是研究 $\mathcal{T}(S)$ 的基本工具。泰希米勒空间还容许整体实解析坐标，称为 Fenchel--Nielsen 坐标。这些坐标依赖于在曲面上选取 $3g-3$ 个互不相同、具有不交代表元的非平凡简单闭曲线同伦类。坐标由各同伦类中测地线的双曲长度以及绕这些测地线的扭转给出。

此外，Thurston 在泰希米勒空间上引入了一个新的黎曼度量：在以某个双曲曲面为代表的点处，两个切向量的内积是曲面上均匀分布的一列闭测地线长度之二阶导数。这个度量只用到曲面的双曲几何，与既有的、基于拟共形理论的泰希米勒度量形成对照，也与 Weil--Petersson 度量形成对照——后者的定义用到 Petersson Hermite 内积，那是在模形式语境中定义、被数论学家使用的。事实上，Wolpert 后来证明 Thurston 的度量与 Weil--Petersson 度量一致。
% [Tutor Note: 维数自查：T(S) 的实维数 6g−6 恰与 Fenchel--Nielsen 坐标的数目一致——3g−3 条简单闭测地线的长度，加上 3g−3 个绕这些测地线的扭转参数，共 6g−6。段末"Thurston 度量 = Weil--Petersson 度量"由 Wolpert 证明；作者强调"泰希米勒度量基于拟共形理论"，正是为了对照 Thurston 度量只用到双曲几何、而 Weil--Petersson 度量来自模形式的 Petersson 内积。]

% [原文 p.xxix]

关于这一极其丰富的现代几何领域的详尽而富有教学意味的综述，我们请读者参阅 Imayoshi 与 Taniguchi 的书 [212]，以及 Hubbard 的书系列 [211] 和由 Papadopoulos 主编的七卷本《泰希米勒理论手册》[295]。泰希米勒空间的各个方面将在本书后面各章中详细研究。

\subsection*{Hitchin 分支}

亏格 $g\ge 2$ 的曲面 $S$ 的泰希米勒空间具有一个更代数的解释：它是基本群 $\pi_1(S)$ 到群 $\operatorname{PSL}_2(\R)$（双曲平面的保向等距群）的不可约、离散且忠实表示构成的空间。这一替代而影响深远的视角归功于 Fricke，可以通过所谓的和乐表示看出。$S$ 上的双曲结构 $(M,f)$ 给出同态 $\rho\colon\pi_1(S)\to\operatorname{PSL}_2(\R)=\operatorname{Isom}^+(\HH^2)$，而由等价双曲结构诱导的表示在 $\operatorname{PSL}_2(\R)$ 中元素共轭的意义下相同，反之亦然。于是
\[
\mathcal{T}(S)\subset\mathcal{R}(S,\operatorname{PSL}_2(\R)):=\Hom^+(\pi_1(S),\operatorname{PSL}_2(\R))/\operatorname{PSL}_2(\R),
\]
即基本群到群 $\operatorname{PSL}_2(\R)$ 的约化表示的模空间。后者称为群 $\operatorname{PSL}_2(\R)$ 的特征簇。事实上，如 Goldman 于 1980 年所证明，$\mathcal{T}(S)$ 是特征簇 $\mathcal{R}(S,\operatorname{PSL}_2(\R))$ 中的一个连通分支。它是完全由离散忠实表示构成的两个连通分支之一；另一个这样的分支是 $\mathcal{T}(\bar S)$，其中 $\bar S$ 表示定向相反的曲面。

从泰希米勒空间作为 $\operatorname{PSL}_2(\R)$-特征簇的连通分支这一代数认识出发，自然产生的问题是：对于秩高于 $\operatorname{SL}_2(\R)$ 的群，能否在特征簇 $\mathcal{R}(S,G)$ 中识别出在曲面上具有特殊几何意义的连通分支？20 世纪对经典泰希米勒空间的多面性研究，使这一方向成为一个充满希望的新研究方向。

Hitchin 把泰希米勒空间与黎曼曲面上自对偶 Yang--Mills 方程的某些解联系起来所带来的新洞见，很快给出了在高阶中实现这一构想的第一批连通分支的例子。即，对伴随分裂实半单李群 $G$，存在唯一的嵌入 $\sigma\colon\SL_2(\R)\to G$，它是到 $\mathfrak{g}$ 的主三维子代数的相伴李群同态，该子代数称为 Kostant 主子代数 $\mathfrak{sl}_2(\R)\subset\mathfrak{g}$。在固定一个离散嵌入 $\iota\colon\pi_1(S)\to\SL_2(\R)$ 之后，Hitchin [199] 于 1992 年用解析工具证明：包含 $\sigma\circ\iota\colon\pi_1(S)\to G$ 的连通分支在拓扑上是一个实维数为 $(2g-2)\dim G$ 的球。特别地，当群为 $G=\operatorname{PSL}_2(\R)$ 时，这个分支恰好就是泰希米勒空间。这些连通分支如今称为 Hitchin 分支。

% [原文 p.xxx]

然而，正如 Hitchin 在其原始工作 [199] 中已指出的，用于实现这些连通分支的解析工具并没有给出关于其几何意义的任何提示。这一点从 2006 年起随着 Fock--Goncharov [129]、Guichard [175] 与 Labourie [242] 的工作开始变得清楚。特别地，位于 Hitchin 分支内的基本群表示都是离散、忠实且 Anosov 的。后者是 Labourie 引入的概念，指的是编码了某种与 Anosov 流相类似的动力结构的表示。此外，人们还得到了 Hitchin 分支中表示的一个重要刻画，它涉及一个连续等变映射的存在性：该映射把 $\R P^1$ 中的正三元组送到与李群 $G$ 相联系的某些旗簇中的正三元组。这一刻画结果表明，它实质上为更整体地处理"如何为高阶李群识别经典泰希米勒空间的推广"这一问题开辟了道路。

% [Tutor Note: 这一段的"Anosov 表示"与"连续等变边界映射"是理解 ch10–ch12 的关键线索：前者由 Labourie 2006 年引入（Anosov 表示至今仍是 Hitchin 分支研究的主线之一），后者正是正性/Θ-正性（ch12 主题）的几何表述——把 RP^1 的循环序与正三元组结构"搬运"到 G 的旗簇上。读 ch12 时可将此处与定义 12.1 前后对照：正三元组在旗簇中的像仍满足横截/正性条件，这是"极大边界映射"在分裂实群情形的类比。]

\subsection*{高阶泰希米勒空间}

高阶泰希米勒理论研究的是位于特征簇 $\mathcal{R}(S,G)$ 的特定子集中的基本群表示的性质，其中 $G$ 是单实群。为理解这些子集的特征，人们使用了来自几何、规范理论、代数几何与动力系统的丰富方法，其中许多方法由模空间 $\mathcal{R}(S,G)$ 的非阿贝尔 Hodge 理论提供。"高阶泰希米勒空间"一词源自 Fock 与 Goncharov 的工作 [129]：他们在一般分裂实半单约化李群的语境下，对 Lusztig [251] 的全正性概念发展了一种更代数几何化的途径，并定义了 $S$ 的基本群到这些群的正表示（positive representation）。在建立若干重要几何结果的过程中，Fock 与 Goncharov 构造出了全部正表示，并证明它们是忠实、离散且正双曲的。

今天，"高阶泰希米勒空间"一词更一般地指约化李群 $G$ 的特征簇 $\mathcal{R}(S,G)$ 的一个连通分支，它完全由像离散的忠实表示构成。

高阶泰希米勒空间的若干本质特征可以追溯到 Thurston 的思想与工作；现代研究中一个非常活跃的领域，就是以恰当的方式把经典泰希米勒空间享有的各种几何特征建立到这一高阶情形中。例如，Thurston 的剪切坐标已由 Fock 与 Goncharov [129] 推广到此情形，即今天所谓的 Fock--Goncharov 坐标。另一个例子涉及用热力学形式主义（Thermodynamic Formalism）在高阶泰希米勒空间上产生映射类群不变的黎曼度量 [67]。

尽管如此，正是上述关于 Hitchin 分支中表示的刻画（涉及把 $\R P^1$ 中正三元组送到 $G$ 的旗簇中正三元组的连续等变映射的存在性）提供了动机，去提出一种统一的
% [原文 p.xxxi]
正性框架，而这一框架事实上可以区分出所有的高阶泰希米勒空间。为寻找这样一种一般刻画，Guichard 与 Wienhard 在 [179] 中引入了曲面群表示的一个正性概念，它统一了此前为特定李群族 $G$ 引入的各种正性概念。这一刻画归结为存在非常特殊的边界映射，把正三元组送到旗簇 $G/P_\Theta$ 中的正三元组，其中 $P_\Theta<G$ 是由一组单正根 $\Theta$ 决定的抛物子群。事实上，Guichard 与 Wienhard 在 [179] 中证明，对这种正性结构出现的李群有完整的分类。此外，对这些群，特征簇中 $\Theta$-正表示的集合是开且闭的，从而是 $\mathcal{R}(S,G)$ 的连通分支 [51]。

% [Tutor Note: 这里的"由单正根子集 Θ 决定的抛物子群 P_Θ"是 ch12 的主要对象（Θ-正性）：当 Θ 取全部单正根时 P_Θ 为 Borel 子群、得到全正性；当 Θ 为空集时得到 Guichard–Wienhard 意义下的"极大"正性（对应 ch11 的极大表示，取 G = Sp_{2n}(R) 时正是 Lagrange 三元组的极大性）。注释 51 指 Guichard–Wienhard 关于 Θ-正性给出开闭分支的论文，可与 ch12 的定理对照阅读。]

谱网络由 Gaiotto、Moore 与 Neitzke 于 2013 年引入 [154]，它们是一类黎曼曲面上的轨道网络，用来计算各种四维 $\mathcal{N}=2$ 量子场论的 BPS 谱，稍后我们将对此加以讨论。然而，这些轨道网络可以被用来在平坦 $G$-联络的模空间上定义局部坐标系，从而特别用于研究高阶泰希米勒空间。确实，谱网络与高阶泰希米勒理论之间最清楚的关系之一是由 Gaiotto、Moore 与 Neitzke 在 [156] 中建立的：他们证明，平坦联络模空间上来自特殊谱网络的 Darboux 坐标与 Fock--Goncharov 坐标一致。此外，Hollands 与 Neitzke 在 [205] 中证明，有两类特殊的谱网络，分别与黎曼曲面的理想三角剖分和裤子分解相联系，它们分别导致曲面上平坦 $\SL_2(\C)$-联络模空间上的 Fock--Goncharov 坐标与复化 Fenchel--Nielsen 坐标。此后，谱网络在研究平坦联络模空间及几何中相关对象方面的效力迅速增长。

\section*{超对称量子场论}

\subsection*{超对称}

1967 年的 Coleman--Mandula 定理 [94] 对可由 $S$-矩阵描述的物理量子场论的对称性施加了严格限制，其中 $S$ 代表"散射"。这里"物理的"指一整套精确的公理，而 $S$-矩阵是一个酉矩阵，它把散射过程中涉及的所谓入态集合映到出态集合。该定理断言：如果这样一个量子场论的对称性构成李群，那么它必定是时空对称性的 Poincar\'e 群与某个内部对称性群的直积。然而，超对称量子场论扩展了对称结构却不违反 Coleman--Mandula 定理，因为它们的对称性构成超李群而非通常的李群。这一更大的对称群使超对称理论得以绕过该定理的限制，同时为统一内部对称性与时空对称性提供了更丰富的框架。物理超对称 QFT 的最一般对称超群受 Haag--Lopuszanski--Sohnius 定理 [182] 的约束。

% [原文 p.xxxii]


超对称是把玻色子与费米子相互联系的时空对称性。在超对称量子场论（QFT）中，无穷小对称性构成超李代数，即 $\Z_2$-分次李代数。除通常的 Poincar\'e 生成元外，四维超庞加莱代数还包含奇生成元 $Q_\alpha$ 与 $\bar Q_{\dot\alpha}$，其中 $\alpha=1,2$、$\dot\alpha=\dot 1,\dot 2$，它们就是超荷（supercharge）。关键之处在于，超荷不与 Lorentz 生成元交换，这意味着它们改变粒子的自旋，从而把费米子映为玻色子、反之亦然。此外，它们满足如下反对易关系：
\[
\{Q_\alpha,\bar Q_{\dot\beta}\}\propto\sigma^\mu_{\alpha\dot\beta}P_\mu,
\]
其中 $P_\mu$ 是时空平移的生成元。这样，超荷可以解释为时空平移的"奇平方根"。四维超庞加莱代数还有扩展版本。在 $\mathcal{N}$-扩展的超庞加莱代数中，有 $\mathcal{N}$ 份奇生成元 $Q^I_\alpha$ 与 $\bar Q^I_{\dot\alpha}$，其中 $I=1,\dots,\mathcal{N}$。如果人们关心的是真正的量子场论——特别是不包含自旋大于 1 的基本粒子的理论——那么超对称的数目 $\mathcal{N}$ 必须小于或等于 4。

1974 年，Wess 与 Zumino 观察到超对称（SUSY）赋予量子场论以极好的重整化性质 [359]，使其从唯象角度看极具吸引力。SUSY 很快作为粒子物理标准模型若干缺陷的一个有希望的解决方案而出现，特别是在处理等级问题（hierarchy problem）、规范相互作用大统一的可能，以及为暗物质组分提供候选者等方面。标准模型的超对称扩展即所谓最小超对称标准模型（MSSM），于 1981 年被提出 [110]。20 世纪 80 年代见证了许多朝超对称唯象实现方向的进展，特别是动力学超对称破缺的发现与研究 [362, 5]。

超对称量子场论卓越的重整化性质被概括为所谓的 SUSY 非重整化定理 [172]。这些定理在 Seiberg 于 20 世纪 90 年代初发展全纯性（holomorphy）概念时起了关键作用 [316, 317, 318, 319]。全纯性断言 SUSY 理论中某些重整化后的量是（裸）场与耦合的解析函数，而非更一般的平滑函数。这样，超对称量子场论与普通量子场论之间的区别就映照出解析几何与光滑几何之间的差别。这一深刻的结构对超对称理论中的量子修正施加强约束，使其比非超对称对应物更易处理。这一全纯原理最早的成功应用之一是对 $\mathcal{N}=1$ 超 QCD（SQCD）的详细研究 [318, 319]，它导致了 Seiberg 对偶的发现 [312]。Seiberg 对偶是一个重要结果，它特别揭示了：规范群为 $\SU(N)$、含 $N+2\le F\le 3N$ 种基本夸克味的 $\mathcal{N}=1$ SQCD，与规范群为 $\SU(F-N)$、含 $F$ 种夸克的修正版本 $\mathcal{N}=1$ SQCD，
% [原文 p.xxxiii]
是同一宏观物理的不同微观描述。这一对偶为超对称规范理论的非微扰行为提供了深刻洞见。

\subsection*{电磁对偶与 BPS 态}

在 Polyakov、't Hooft 与 Julia--Zee 工作的基础上，Bogomol'nyi 于 1975 年在 [59] 中导出了 Georgi--Glashow 模型中作为一个孤子得到的双子质量的经典不等式：
\[
M_{q_{\mathrm{el}},q_{\mathrm{mag}}}\ge\alpha\sqrt{q_{\mathrm{el}}^2+q_{\mathrm{mag}}^2},
\]
其中 $q_{\mathrm{el}}$ 表示双子的电荷、$q_{\mathrm{mag}}$ 表示其磁荷，$\alpha$ 是一个常数。等号只在底层非交换规范理论所谓的 Bogomol'nyi--Prasad--Sommerfield（BPS）极限下成立 [306]，其中使 Bogomol'nyi 界饱和的双子称为 BPS 态。所有这些结果都是纯经典的，并没有根本理由说明它们在引入量子效应后仍然成立。具体地说，双子的质量原则上可能被量子修正改变。

在具有扩展超对称的量子场论中，特别是对 $\mathcal{N}\ge 2$，超庞加莱代数含有称为中心荷（central charge）的中心生成元。对 $\mathcal{N}=2$，中心荷取复数的形式。具有 $\mathcal{N}=2$ 超对称的超 Yang--Mills 理论，是 Georgi--Glashow 模型在 Bogomol'nyi--Prasad--Sommerfield（BPS）极限下包含费米子的特例推广。1978 年，Witten 与 Olive 证明：这类理论中双子的中心荷从它的低能电荷与磁荷获得贡献 [369]。他们还确立了 BPS 态对应于 $\mathcal{N}=2$ 扩展超庞加莱代数的特定表示，即所谓短表示。更精确地说，在具有 $\mathcal{N}=2$ 超对称的理论中，任何超多重态的质量都以其中心荷 $\mathcal{Z}$ 的绝对值为下界：
\[
M\ge|\mathcal{Z}|.
\]
当不等式严格成立时，该超多重态被归为超对称代数的长表示，此时超荷 $Q^I_\alpha$ 与 $\bar Q^I_{\dot\alpha}$ 中没有一个被平凡地实现。反之，当不等式饱和时，有一半超荷在该表示上作为零算子作用，这意味着短超多重态的维数是长超多重态的一半。由于量子修正无法弥补这一差别，在 $\mathcal{N}=2$ 量子场论中，BPS 态受表示论保护而不受量子修正影响。

与自由的无质量 Maxwell 理论不同，量子 Yang--Mills 理论中的规范耦合会重整化，这使非交换量子 Yang--Mills 理论中 Montonen--Olive 对偶的理解与表述变得复杂。然而，在具有 $\mathcal{N}=2$ 超对称的超 Yang--Mills（SYM）理论中，由于 Witten 与 Olive 提供的洞见，情况显著改善。在这些理论中，BPS 态作为量子态出现，同时也可以经典地分析。

% [Tutor Note: 原文在 p.xxxii（PDF p030）把中心荷印作花体 𝒵，而在 p.xxxix（PDF p037）与后文又印作正体 Z（如 Z_{γ^{ij}}）。译文统一用 $\mathcal{Z}$ 表示中心荷，读者对照原文时请注意这是同一记号；BPS 不等式 M ≥ |𝒵| 以及短/长多重态的论证，正是"BPS 态受表示论保护"这一全书反复使用的核心论据（见 ch22）。]

% [原文 p.xxxiv]
虽然大多数 $\mathcal{N}=2$ SYM 理论受重整化影响，但 $\mathcal{N}=4$ SYM 理论——$\mathcal{N}=2$ 理论的特例——不受影响。这使 $\mathcal{N}=4$ SYM 理论成为陈述与研究 Montonen--Olive 对偶的理想框架，该对偶最初于 1979 年被猜想 [292]。随后的发展提供了支持 Montonen--Olive 对偶实现的充分证据，在此语境中通常称为 S 对偶。如今人们普遍接受，S 对偶确实在 $\mathcal{N}=4$ SYM 理论中实现，从而在量子场论领域与几何 Langlands 纲领之间建立了引人注目的联系 [226]。

\subsection*{Seiberg--Witten 理论}

在 20 世纪 90 年代初全纯性的发现以及对 $\mathcal{N}=1$ SQCD 理论性质的后续洞见的基础上，Seiberg 与 Witten 于 1994 年证明：四维 $\mathcal{N}=2$ 超对称施加的约束如此之强，以致可以精确求解其低能 Coulomb 动力学。这首先在纯 $\mathcal{N}=2$ $\SU(2)$ SYM [313] 中处理，随后在含 $1\le F\le 4$ 种夸克味的 $\mathcal{N}=2$ $\SU(2)$ SQCD [314] 中处理，之后又推广到具有各种规范群的更一般的 $4d$ $\mathcal{N}=2$ SYM 与 SQCD 理论。

超对称量子场论常常呈现出一族连续的不等价量子真空，称为超对称真空的模空间。对 $\mathcal{N}=2$ SYM 理论尤其如此。在这些理论中，模空间含有一个值得注意的子空间，称为库仑分支（Coulomb branch），它由低能物理是 Maxwell 理论的 $\mathcal{N}=2$ 超对称推广的那些真空组成。在数学上，库仑分支是一个微分同胚于 $\C^r$ 的特殊凯勒流形，这意味着存在一组局部全纯函数，记作 $a_i$ 与 $a^D_i$（$i=1,\dots,r$），称为特殊坐标。这些函数概括了低能有效交换规范理论以及低能双子的中心荷，从而是 BPS 态的质量。Seiberg 与 Witten 的结果给出了这些特殊坐标的精确表达式，从而完全求解了这类理论的低能动力学。

具体地说，特殊坐标 $a_i$ 与 $a^D_i$ 是定义在某个辅助黎曼曲面 $\Sigma$（称为 Seiberg--Witten 曲线）上的一个亚纯 1-形式的周期积分。这条曲线依赖于库仑分支内的真空。Seiberg--Witten 曲线被组织为库仑分支上的一个全纯纤维化，而该纤维化完全由 $a_i$ 与 $a^D_i$ 的奇异行为刻画，后者由物理考虑确定。Seiberg 与 Witten 所发展的技巧很快导致了一批引人注目的新 $4d$ $\mathcal{N}=2$ 超共形场论（SCFT）的发现，其中包括 Argyres--Douglas 理论 [24]（1995）与 Minahan--Nemeschansky 理论 [266, 267]（1996）。

% [原文 p.xxxv]

Seiberg--Witten 理论的一个重要成果是理解了：即使在规范耦合会重整化的情形，$4d$ $\mathcal{N}=2$ SYM 理论中也存在 S 对偶的概念，尽管这一实现比 $4d$ $\mathcal{N}=4$ SYM 理论中的情形更为复杂。与 $4d$ $\mathcal{N}=4$ SYM 理论不同，一般 $4d$ $\mathcal{N}=2$ 理论中 BPS 态的谱在库仑分支上由度量奇点引起的单值变换下并不保持不变。然而，一个称为穿墙（wall-crossing）的现象确保这种非不变性并不与这些理论中 S 对偶的实现相矛盾。库仑分支上存在余维 1 的边际稳定墙（walls of marginal stability），BPS 态可以穿过这些墙而衰变。这些墙穿过库仑分支上的奇点，这表明当人们绕这样的奇点一圈时，一个 BPS 态在穿墙时可能转变为两个态。

1995 年，Donagi 与 Witten 证明：一般 $4d$ $\mathcal{N}=2$ 理论在库仑分支上的低能物理结构定义了一个复完全可积系统 [114]。这一框架涵盖了此前用 Seiberg 与 Witten 方法研究过的所有理论，以及新的理论即 $4d$ $\mathcal{N}=2^*$ 理论。与这些 $4d$ $\mathcal{N}=2$ 量子场论相联系的复完全可积系统，是 Hitchin 在 [198, 197] 中引入的可积系统及其到非闭黎曼曲面推广的特例。

超对称尚未在粒子对撞机实验中被观测到，这使人们对其在 MSSM 等模型中的具体实现产生怀疑，并使标准模型中未解决问题的自然解决方案为何这一问题保持开放。然而，即使 SUSY 并未完全如 MSSM 中那样实现，它仍然是处理高能物理中许多未解问题的一个极具吸引力的框架。在弦论界，SUSY 常被视为弦论的自然预言，尽管这并不必然意味着它应当出现在当前对撞机可及的能量标度上。这种不确定性提出了重要的科学与哲学问题。然而，无论最终答案如何，超对称无疑在推进量子场论形式体系的研究中起了关键作用，即使这一追求更偏向数学领域。上面强调的结果表明了它对理解量子场论的深刻影响。

\section*{谱网络}

\subsection*{弦论与膜}

弦论最早于 20 世纪 70 年代初以对偶共振模型（dual resonance model）之名发展起来，最初是为了描述强核相互作用。这一时期的一个重要产物是 1971 年引入的超弦 Ramond--Neveu--Schwarz 形式体系。然而，随着量子色动力学被公认为
% [原文 p.xxxvi]
强相互作用更准确的描述，对偶共振模型被搁置了。

1974 年弦论复兴：人们发现闭弦的一种振动模式具备可能描述引力子的所需性质，从而使弦论成为统一理论的一个候选者。在其微扰框架中，弦论把弦描述为基本对象，而非点状粒子。在大于典型弦长的尺度上，弦有效地表现为点状粒子，其质量与自旋由弦的具体振动模式决定。

20 世纪 70 年代末与 80 年代弦论发展的重要里程碑包括：1977 年 Gliozzi--Scherck--Olive 投影的发现，1984 年通过 Green--Schwarz 机制实现的反常抵消，以及 1985 年 Gross、Harvey、Martinec 与 Rohm 建立的杂化弦理论——它使十维中已知超弦理论的数目达到五个：I 型、IIA 型与 IIB 型、HO 型与 HE 型。同样在 1985 年，Calabi--Yau 流形对于紧化额外维度的重要性得到认识（Candelas、Horowitz、Strominger 与 Witten，1985）。这对 20 世纪 80 年代与 90 年代之交镜像对称的发现起了关键作用（Dixon，1988；Lerche、Vafa 与 Warner，1989；Candelas、de la Ossa 与 Parks，1991）。镜像对称的研究极大地得益于 Witten 于 1990 年引入的拓扑弦论，这尤其导致了 Kontsevich 的同调镜像对称猜想（1994）。
% [Tutor Note: 原文把 Scherk 拼作 "Scherck"（正确拼法为 Gliozzi--Scherk--Olive，即 1977 年提出的 GSO 投影），此处照原文译作 Scherck，正确拼法见文末文献表。GSO 投影的作用是截断弦谱中的快子并保持时空超对称，是 1984 年 Green--Schwarz 反常抵消之前的关键一步。]

所有十维超弦理论的低能极限就是所谓的十维超引力理论，也称为 I 型、IIA、IIB、HO 与 HE 型。20 世纪 90 年代初，人们越来越清楚：全部五种超引力与超弦理论版本都由对偶——特别是 S 对偶与 T 对偶——相联系，这表明它们应被视为同一个底层理论的不同描述，而非彼此不同的理论。1995 年，Witten 提出 M 理论作为 IIA 型超弦的强耦合极限。与十维弦论不同，M 理论是一个十一维的膜理论。它的低能极限是十一维超引力，后者在某种意义上在超引力理论中是普适的。

1995 年的另一个突破是 Polchinski 证明超弦理论需要纳入弦之外的额外扩展对象：D 膜 [303]。例如，IIA 型超弦容许 D0-、D2-、D4-、D6- 与 D8-膜，它们分别对应于时空中 1、3、5、7、9 维的对象。D 膜是开弦可以终结的所在。在低能下，终结于 D 膜上的开弦的模式在膜上诱导出超 Yang--Mills 理论，并保持 IIA 型超弦中总超荷的一半——具体说是 16 个超荷。换言之，超对称 QFT 在超弦理论（以及 M 理论）中无处不在。除 D 膜与基本弦（常称为 NS1 膜）外，IIA 型超弦还容许所谓的 NS5 膜。

Polchinski 的突破之后不久，人们就认识到：基本弦能终结于 D 膜，这一事实经由一系列对偶意味着膜常常也能终结于其他膜上 [337, 347]。这为
% [原文 p.xxxvii]
研究膜构型的动力学打开了大门，而这些构型在其世界体积（即它们在时空中扫出的体积）上承载超对称 QFT。超对称 Yang--Mills 理论的许多非平凡动力学都可以通过这样的膜构型有效地探索，正如 Hanany 与 Witten 的开创性工作 [186] 等所展示的那样。

D 膜在全息对偶的发现中也起了关键作用：在全息对偶下，$d$ 维量子场论等价于 $(d+1)$ 维量子引力理论（或更一般地，广义协变理论）。全息对偶最突出的例子是 Maldacena 于 1997 年发现的 AdS/CFT 对应 [254]，它把 $S^4$ 上四维 $\mathcal{N}=4$ $\SU(N)$ SYM 理论与 $AdS_5\times S^5$ 上具有 $N$ 单位 Ramond--Ramond 通量穿过 5-球面 $S^5$ 的 IIB 型超弦理论联系起来。关于 AdS/CFT 对应许多方面的早期综述可参见 [6]。此后，全息对偶被显著推广，以纳入超对称更少以及非共形的量子场论。在基本层面上，全息对偶为特定背景上的弦论提供了非微扰定义。

与上一节最后一段相呼应，弦论究竟必须以何种精确方式实现（如果它能实现的话）才能描述我们宇宙的物理，至今仍难以捉摸，这提出了重要的科学与哲学问题。然而，大量证据表明弦论是一个自洽的量子引力理论，这为其作为处理理论高能物理中深层问题（特别是量子引力的本性）的宝贵工具而加以研究提供了理由。此外——而这正是本书将要关注的主要方面——弦论已被证明是研究超对称量子场论的有力工具，增进了我们对一般量子场论框架的理解，并阐明了基础数学与物理之间丰富的相互作用。

\subsection*{来自 M 理论的 Seiberg--Witten 曲线}

一大类四维 $\mathcal{N}=2$ 超 Yang--Mills（SYM）理论可以通过 IIA 型超弦理论中拉伸于 NS5 膜之间的 D4 膜构型来研究，其中所关心的 SYM 理论作为 D4 膜上的世界体积理论出现。用膜构型研究这些 SYM 理论低能物理的一个重要困难在于：在所考虑的构型中 D4 膜终结于 NS5 膜上，而 D4 膜在 NS5 膜上的端点本性一直难以完全理解。1997 年，Witten 证明当把这些膜构型提升到 M 理论时，这一问题就消失了 [366]。通过同时增大 IIA 型弦耦合并对时空中某些方向作重标度，人们可以达到强耦合区域，在那里 M 理论提供更恰当的描述，而不会改变所研究的 $4d$ $\mathcal{N}=2$ 理论。在 M 理论中，系统一般对应于一个包裹光滑黎曼曲面 $\Sigma$ 的单个 M5 膜，而 $\Sigma$ 就是该 $4d$ $\mathcal{N}=2$ 理论的 Seiberg--Witten 曲线。于是，这一基于膜的途径提供了构造
% [原文 p.xxxviii]
无穷族 $4d$ $\mathcal{N}=2$ SYM 理论的 Seiberg--Witten 曲线的统一框架。这些理论可以用所谓特殊酉规范群与味群的 $\mathcal{N}=2$ 线性箭图来描述。编码任何这样一个 $4d$ $\mathcal{N}=2$ SYM 理论低能物理的代数完全可积系统，都是带穿孔球面或环面上的 Hitchin 可积系统。

\subsection*{S 类理论}

人们可以如下分析涉及特殊酉规范群与味群的一般 $\mathcal{N}=2$ 共形线性箭图的 S 对偶性质：假定除一个规范群外其余规范群都极弱耦合，然后每次对一个规范因子作对偶化。这一途径自然地导致对由广义箭图 [146] 描述的四维 $\mathcal{N}=2$ 理论的研究，这些理论不仅涵盖此前研究过的情形 [366]，还包含更多一般情形。这些广义箭图对应于 S 类中 $4d$ $\mathcal{N}=2$ 理论的特定描述。一个 S 类理论由以下数据确定：一个带穿孔的黎曼曲面 $C$，一个 $A$、$D$ 或 $E$ 型的紧单李代数 $\mathfrak{g}$，以及 $C$ 的穿孔处的某些约束 $D$（其中 $D$ 代表"缺陷"，defects）。S 类中的理论通常记作 $\mathcal{S}[C,\mathfrak{g},D]$。

存在具有 $\mathcal{N}=(2,0)$ 超对称的非凡六维超共形场论（SCFT），它们由 $A$、$D$ 或 $E$ 型的紧单李代数 $\mathfrak{g}$ 分类。这些理论容许一组特定的余维 2 缺陷。S 类理论可以统一地用这些六维 SCFT 来描述：理论 $\mathcal{S}[C,\mathfrak{g},D]$ 对应于把 $\mathfrak{g}$ 型的六维 SCFT 在曲面 $C$ 上紧化（带有部分拓扑扭转）、并在 $C$ 的穿孔处插入余维 2 缺陷所得到的 $4d$ $\mathcal{N}=2$ 理论。六维 $\mathcal{N}=(2,0)$ 理论的模空间含有库仑分支，在该分支上理论在低能下约化为交换的六维 $\mathcal{N}=(2,0)$ 理论。这个库仑分支反过来又使人们可以分析由它导出的 S 类 $4d$ $\mathcal{N}=2$ 理论的库仑分支 [146, 155]。

对 $\mathfrak{g}=\mathfrak{su}(N)$，六维 $\mathcal{N}=(2,0)$ 理论描述一叠 $N$ 个 M5 膜上的世界体积理论。库仑分支对应于在与原膜叠横向的方向上分离这些 M5 膜。因此，理论 $\mathcal{S}[C,\mathfrak{su}(N),D]$ 由把一叠 $N$ 个 M5 膜包裹在 $C$ 上、并以由缺陷 $D$ 确定的边界条件来定义。这个四维理论的库仑分支对应于 M5 膜叠在 $C$ 的全纯余切丛中的可归一化形变。这些形变通常定义 $C$ 的一个 $N$ 次全纯分歧覆叠，它被等同于 $\mathcal{S}[C,\mathfrak{su}(N),D]$ 在其库仑分支上某个特定点处的 Seiberg--Witten 曲线。

支配 S 类理论库仑分支上低能物理的可积系统是 Hitchin 系统。与四维情形中库仑分支只对应于 Hitchin 纤维化的底不同，当 S 类理论在圆周 $S^1$ 上紧化时，整个 Hitchin 系统——它是超凯勒的——作为紧化理论的库仑分支出现 [322]。为从六维导出这一点，可以先在 $S^1$ 上紧化六维 $\mathcal{N}=(2,0)$ 理论，这
% [原文 p.xxxix]
给出一个五维 SYM 理论。再在黎曼曲面 $C$ 上（带部分拓扑扭转）作进一步紧化，就得到 $\R^3\times S^1$ 上的 S 类理论。五维理论在 $C$ 上紧化后的 BPS 方程恰好就是 Hitchin 方程 [155]。

\subsection*{超凯勒度量与穿墙现象}

一般地，在半径为 $R$ 的圆周 $S^1$ 上紧化的 $4d$ $\mathcal{N}=2$ 理论具有一个带光滑超凯勒度量 $g_R$ 的超凯勒库仑分支。当 $R\to\infty$ 时，该度量 $g_R$ 趋近于某个特定的超凯勒"半平坦"度量 $g^{\mathrm{sf}}$。虽然 $g^{\mathrm{sf}}$ 可以显式写出，但它在余维 2 的流形上呈现奇点。量子修正在有限的 $R$ 处出现并使这些奇点光滑化，从而给出非奇异的度量 $g_R$。这些量子修正只依赖于该 $4d$ $\mathcal{N}=2$ 理论的 BPS 态。关键在于，$g_R$ 保持光滑这一要求对 BPS 谱的穿墙行为施加约束 [151]：它们必须满足 Kontsevich--Soibelman 穿墙公式，该公式在 [237] 中于不同语境下导出。超凯勒度量 $g_R$ 的（扭量）构造涉及一组关键的对象，即所谓的 Darboux 坐标 $\mathcal{X}^{\mathrm{RH}}_\gamma$。

在 S 类 $4d$ $\mathcal{N}=2$ 理论这一特定情形，三维库仑分支对应于一个 Hitchin 系统，即黎曼曲面 $C$ 上平坦联络的模空间。在此语境中，Darboux 坐标 $\mathcal{X}^{\mathrm{RH}}_\gamma$ 可由 Fock 与 Goncharov [129] 引入的平坦联络模空间上的坐标导出。Gaiotto、Moore 与 Neitzke 的奠基性工作 [155] 标志着在重新联系高阶泰希米勒理论——具体说是 Hitchin 分支——与 S 类 $4d$ $\mathcal{N}=2$ 理论研究方面迈出的重要一步。

\subsection*{WKB 网络}

在 $A$ 型的 S 类理论中（它们来自把 $N$ 个 M5 膜在黎曼曲面 $C$ 上紧化），库仑分支上一般点处的 BPS 态——该点对应于一个 $N$ 次分歧覆叠 $\Sigma\to C$——由拉伸于覆叠各叶之间的 M2 膜表示。这样的 M2 膜到 $C$ 上的投影构成一条曲线。局部地，覆叠的各叶可以用 $1,\dots,N$ 来标记，从而可以用 $\gamma^{ij}$ 来标记这些曲线。对每个 $i=1,\dots,N$，设 $\lambda_i$ 表示 Seiberg--Witten 微分 $\lambda$ 到第 $i$ 叶的限制。与 $\gamma^{ij}$ 相联系的中心荷 $Z_{\gamma^{ij}}$ 与质量 $M_{\gamma^{ij}}$ 由下式给出：
\[
Z_{\gamma^{ij}}=\frac{1}{\pi}\int_{\gamma^{ij}}\lambda_i-\lambda_j,
\qquad\text{以及}\qquad
M_{\gamma^{ij}}=\frac{1}{\pi}\int_{\gamma^{ij}}|\lambda_i-\lambda_j|.
\]
BPS 不等式 $M_{\gamma^{ij}}\ge|Z_{\gamma^{ij}}|$ 自动成立。该不等式饱和当且仅当 $\lambda_i-\lambda_j$ 的辐角沿 $\gamma^{ij}$ 保持不变。因此，四维理论中的 BPS 态由 $C$ 上满足下式的曲线描述：

% [原文 p.xl]

\[
\lambda_{ij}\cdot\partial_s\in \ee^{\ii\vartheta}\R_+,
\]
其中 $s$ 是 $\gamma^{ij}$ 上的坐标，$\vartheta\in[0,2\pi)$。对任意 $\vartheta$ 值，这个微分方程都定义 $C$ 的一个奇异叶状结构（在 $N=2$ 的情形）。当 $N\ge 2$ 时，$C$ 上可以有 $ij$-曲线的交汇点。这种把 BPS 态表示为曲面上曲线的做法，最早是在几何工程化的 $4d$ $\mathcal{N}=2$ 理论语境中研究的 [233, 328]。
% [Tutor Note: 记号提醒：λ_ij := λ_i − λ_j，其中 λ_i、λ_j 分别是 Seiberg--Witten 微分 λ 在第 i、j 叶上的限制，∂_s 是曲线 γ^{ij} 的切向。BPS 条件 M = |Z| 等价于 arg(λ_i − λ_j) 沿曲线为常数 ϑ，故有 (λ_i − λ_j)·∂_s ∈ e^{iϑ}R_+。N = 2 时该方程给出 C 的奇异叶状结构；N ≥ 2 时不同的 ij-曲线可在覆叠的分支点处交汇，这正是 WKB 网络/谱网络"三叉点"的几何来源（图 0.1 中的橙色叉号即分支点）。]

与 $C$ 上满足这个微分方程的一族曲线（一个"网"）相联系的四维 BPS 态具有有限质量，当且仅当该网的总长度有限。这自然地引向考虑 $C$ 上源自覆叠 $\Sigma\to C$ 的分支点的特定网，在分支点处三条积分曲线相交。这些网定义了给定角 $\vartheta$ 下 $C$ 上的 WKB 网络 $\mathcal{W}_\vartheta$ [155]（因 WKB 近似而得名）。对一般的 $\vartheta$ 值，不存在有限长度的网。然而，随着 $\vartheta$ 变化，WKB 网络演化并在某些特定的 $\vartheta$ 值处发生突变，此时有限网（等价地，真正的 BPS 态）出现。通过让 $\vartheta$ 从 $0$ 连续变到 $\pi$，人们可以描述理论在其库仑分支上给定点处的全部 BPS 态。纯 $4d$ $\mathcal{N}=2$ $\SU(2)$ 超 Yang--Mills 理论（其中 $C=\C P^1$）在库仑分支原点处、不同 $\vartheta$ 值下的 WKB 网络例子见图 0.1。在 $\vartheta=0$ 与 $\vartheta=\pi/2$ 处，$C$ 上出现有限的 BPS 曲线，对应该理论在库仑分支原点处的两个 BPS 态——即单极子与荷为 $(2,1)$ 的双子。

\subsection*{谱网络}

谱网络是本书的中心主题，它于 [154] 中被正式引入。它们建立在带标架 BPS 态的概念之上，后者是在四维理论语境中引入的 [153]，后来又被推广到耦合的 $2d$-$4d$ 系统 [152]。

S 类理论容许典范的半 BPS 面缺陷 $S_z$，它们由点 $z\in C$ 参数化 [147, 152, 10]。当 $z\in C$ 是对应于该理论库仑分支上某个固定点的覆叠 $\Sigma\to C$ 的一般点时，缺陷 $S_z$ 容许 $N$ 个有质量真空，它们对应于 $z$ 在 $\Sigma$ 中的 $N$ 个原像。在存在面缺陷 $S_z$ 时，理论支持称为孤子（soliton）的 BPS 粒子，它们束缚于该缺陷并在 $S_z$ 的不同真空之间插值。因此，孤子的荷 $a$（大致上）是 $\Sigma$ 中 $z$ 的两个不同提升之间道路的同调类（相对于其边界）。荷为 $a$ 的孤子与一个记作 $Z_a(z)$ 的中心荷相联系。几何上，孤子被表示为 $C$ 上有限的弦网，与前面讨论的类似，只是其中一条弦必须终结于 $z$。孤子的简并度由 BPS 指标 $\mu(a)$ 捕捉。来自 WKB 网络的 BPS 孤子在讲义 [63] 与 [279] 中有详细研究。

% [原文 p.xli]

\Fig[0.95\linewidth]{figures/fig_0_1.png}{0.1}{在库仑分支 $C$ 的原点处、某些 $\vartheta$ 值下的 BPS 曲线。橙色叉号表示分歧覆叠 $\Sigma\to C$ 的分支点，蓝色圆点表示 $C=\C P^1$ 上的一个穿孔——此处用球极投影把它画成一个平面（在无穷远处还有一个穿孔）。}

对参数 $\vartheta\in\R$，$C$ 上的谱网络定义为满足下述条件的点 $z\in C$ 的集合：存在 $\Sigma$ 中 $z$ 的两个不同提升之间的同调类 $a$（相对于其边界），使得 $Z_a(z)/\ee^{\ii\vartheta}\in\R_-$ 且 $\mu(a)\ne 0$。

给定两个点 $z_1,z_2\in C$，以及一条从 $z_1$ 到 $z_2$ 的道路 $\wp$ 与一个实参数 $\vartheta$，可以定义面缺陷 $S_{z_1}$ 与 $S_{z_2}$ 之间的一个超对称界面 $L_{\wp,\vartheta}$，它保持两个超荷。$2d$-$4d$ 带标架 BPS 态对应于 $\vartheta$ 取一般值时界面 $L_{\wp,\vartheta}$ 的不同真空，它们在 $\Sigma$ 中 $z_1$ 的一个提升与 $z_2$ 的一个提升之间插值。与孤子情形一样，这些 $2d$-$4d$ 带标架 BPS 态的简并度被编码在一个记作 $\Omega'(L_{\wp,\vartheta},a)$ 的指标中，其中 $a$（大致上）是从 $z_1$ 的一个提升到 $z_2$ 的一个提升的道路的同调类。带标架的 $2d$-$4d$ 指标被汇集在一个记作 $F(\wp,\vartheta)$ 的生成函数中，它是从 $z_1$ 的提升到 $z_2$ 的提升的全部道路同调类按带标架 $2d$-$4d$ 指标加权所得的线性组合。对任意从 $z_1$ 到 $z_2$ 的道路 $\wp_1$ 与从 $z_2$ 到 $z_3$ 的道路 $\wp_2$，生成函数 $F(\wp_1,\vartheta)$ 与 $F(\wp_2,\vartheta)$ 可以复合，其复合来自 $\Sigma$ 上路径代数中的连接乘积。函数 $F(\wp,\vartheta)$ 关键地依赖于道路 $\wp$ 如何与谱网络 $\mathcal{W}_\vartheta\subset C$ 相交。

物理上，超对称界面 $L_{\wp,\vartheta}$ 预期只依赖于道路 $\wp$ 的同伦类。这一同伦不变性质对 $F(\wp,\vartheta)$ 与孤子简并度 $\mu(a)$ 都施加严格约束。此外，虽然 $F(\wp,\vartheta)$ 在 $\vartheta$ 取一般值时是局部常值的，但存在特定的 $\vartheta$ 值
% [原文 p.xlii]
使 $F(\wp,\vartheta)$ 发生跳跃。这些跳跃恰好发生在四维理论具有中心荷辐角为 $\ee^{\ii\vartheta}$ 的 BPS 态的那些 $\vartheta$ 值处。此外，四维 BPS 态的简并度可以通过分析函数 $F(\wp,\vartheta)$ 的行为推断出来。因此，谱网络提供了一个框架，用于在 Lie 代数 $\mathfrak{g}$ 的秩大于 1 的情形计算 S 类理论的四维 BPS 谱，极大地推广了用 [155] 中引入的 WKB 网络所得的结果。

生成函数 $F(\wp,\vartheta)$ 以同伦不变的方式把 $\Sigma$ 上道路的组合与 $C$ 上的道路 $\wp$ 相联系，这一事实是关键特征。该性质使这些函数可以被用作从 $\Sigma$ 上的交换平坦联络构造 $C$ 上非阿贝尔平坦联络的工具。这一过程称为非阿贝尔化（non-abelianization），而逆过程称为阿贝尔化（abelianization）[205]。在 [154] 中，这一性质被作为指导原则，用来定义谱网络的拓扑概念——它可以看作与一个未必全纯的分歧覆叠 $\Sigma\to C$ 相联系的 $C$ 上的图。它也是谱网络作为研究高阶泰希米勒空间的工具而日益受到关注的主要原因。虽然谱网络最初在 [154] 中被引入，是为了推广 Fock--Goncharov 坐标（$A$ 型的，以及稍作修改的 $X$ 型的），但如今它们已被证明涵盖了更多具有几何与物理意义的话题。

在结束这篇介绍谱网络的简短概述时，我们提到：绘制谱网络的一个非常有效的工具是 P. Longhi、C. Park 与 A. Neitzke 开发的程序"Loom"，它在 GitHub 上可免费获取 [250]。
本书的封面图以及图 0.1 与图 22.17 都是用 Loom 得到的。

% [Tutor Note: 图 0.1 由 Loom 生成，8 个分图对应 ϑ = 0.081π、0.333π、0.455π、0.5π、0.556π、0.677π、0.919π 与 π（≈0）；其中 ϑ = 0 与 ϑ = π/2 处出现有限长度的 BPS 曲线，对应该纯 SU(2) 理论在库仑分支原点处的两个 BPS 态（电荷 (2,1) 的单极子与双子），这正是 ch22 要定量计算的对象。图上 8 个分图按循环顺序排列，ϑ = π 与 ϑ = 0 重合，体现了谱网络的 2π 周期性。]
